%-------------------------
% Resume in Latex
% Author : Jake Gutierrez
% Based off of: https://github.com/sb2nov/resume
% License : MIT
%------------------------

% This is the editor's port of Jake's Resume, from https://github.com/jakegut/resume, which is
% based on Sourabh Bajaj's resume, https://github.com/sb2nov/resume, so the template derives from
% both. Each is under the MIT licence below, whose permission notice is the same for the two, and
% the copyright lines are theirs. The port changes only what the editor's compile loop needs.
% geometry sets the page size and margins, in place of fullpage and the lines that adjusted its
% margins, so that the page size and margins controls are ordinary compiles. hyperref takes the
% dvipdfmx driver, because the loop writes DVI and the editor makes the PDF from it with
% xdvipdfmx, and makes no bookmarks, which a one-page resume has no use for and which cost time
% that grows with the square of a section title's length. The fonts are Latin Modern in T1, where
% upstream has Computer Modern in OT1: the editor draws the page from Type 1 outlines, and TeX
% Live has none for Computer Modern in T1, while Latin Modern is the same design and has them;
% and in OT1 LaTeX cannot set letters such as ą or Đ, or guillemets, at all. article sets its
% sizes once, as it loads, and only at 10, 11 or 12 points, while the font size control takes half
% points from 9 to 12 as an ordinary compile, so that one format serves every size: the port's
% \rtfontsize sets article's 11-point sizes and the spacing that goes with them, as its size11.clo
% has them, scaled by the size asked for over 11 points and rounded to a hundredth of a point, so
% that every size, the name's, the titles' and the text's, and the space between lines, grows with
% the control, and at 11 points, the class's size and the control's default, the page is
% upstream's to the byte. The port names the accent colour, rtaccent, and sets the section titles
% and the rules under them in it, where upstream set the rules in black, so that the accent colour
% control can set it; its default is black, which keeps upstream's look. xcolor takes the place of
% color, because the control gives the colour as six hex digits, which only xcolor reads. And the
% file ends where the body would begin, because the editor generates the body from the resume.
% Everything else is upstream's.
%
% Copyright (c) 2026 Sourabh Bajaj
% Copyright (c) 2020 Jake Gutierrez
%
% Permission is hereby granted, free of charge, to any person obtaining a copy
% of this software and associated documentation files (the "Software"), to deal
% in the Software without restriction, including without limitation the rights
% to use, copy, modify, merge, publish, distribute, sublicense, and/or sell
% copies of the Software, and to permit persons to whom the Software is
% furnished to do so, subject to the following conditions:
%
% The above copyright notice and this permission notice shall be included in all
% copies or substantial portions of the Software.
%
% THE SOFTWARE IS PROVIDED "AS IS", WITHOUT WARRANTY OF ANY KIND, EXPRESS OR
% IMPLIED, INCLUDING BUT NOT LIMITED TO THE WARRANTIES OF MERCHANTABILITY,
% FITNESS FOR A PARTICULAR PURPOSE AND NONINFRINGEMENT. IN NO EVENT SHALL THE
% AUTHORS OR COPYRIGHT HOLDERS BE LIABLE FOR ANY CLAIM, DAMAGES OR OTHER
% LIABILITY, WHETHER IN AN ACTION OF CONTRACT, TORT OR OTHERWISE, ARISING FROM,
% OUT OF OR IN CONNECTION WITH THE SOFTWARE OR THE USE OR OTHER DEALINGS IN THE
% SOFTWARE.

\documentclass[letterpaper,11pt]{article}

% Latin Modern in T1, in place of Computer Modern in OT1, as the note at the top says.
\usepackage[T1]{fontenc}
\usepackage{lmodern}
% The font size control's sizes, as the note at the top says. \rt@size is the size in scaled points,
% 11 points until \rtfontsize sets another, and \rt@pt{<points>} is a length of size11.clo's scaled
% by it over 11 points and rounded to a hundredth of a point, TeX arithmetic throughout. A length
% to two decimals rounds back to what TeX reads for it, so at 11 points every length is article's.
\makeatletter
\newcommand*\rt@size{720896}
\newcommand*\rt@pt[1]{%
  \the\dimexpr\numexpr\dimexpr#1\p@*\rt@size/720896\relax*100/65536\relax\p@/100\relax}
% size11.clo's sizes, with the display skips and the lists' spacing it sets with them, each length
% through \rt@pt.
\DeclareRobustCommand\normalsize{%
  \@setfontsize\normalsize{\rt@pt{10.95}}{\rt@pt{13.6}}%
  \abovedisplayskip \rt@pt{11}\@plus\rt@pt{3}\@minus\rt@pt{6}%
  \abovedisplayshortskip \z@ \@plus\rt@pt{3}%
  \belowdisplayshortskip \rt@pt{6.5}\@plus\rt@pt{3.5}\@minus\rt@pt{3}%
  \belowdisplayskip \abovedisplayskip
  \let\@listi\@listI}
\DeclareRobustCommand\small{%
  \@setfontsize\small{\rt@pt{10}}{\rt@pt{12}}%
  \abovedisplayskip \rt@pt{10}\@plus\rt@pt{2}\@minus\rt@pt{5}%
  \abovedisplayshortskip \z@ \@plus\rt@pt{3}%
  \belowdisplayshortskip \rt@pt{6}\@plus\rt@pt{3}\@minus\rt@pt{3}%
  \def\@listi{\leftmargin\leftmargini
    \topsep \rt@pt{6}\@plus\rt@pt{2}\@minus\rt@pt{2}%
    \parsep \rt@pt{3}\@plus\rt@pt{2}\@minus\rt@pt{1}%
    \itemsep \parsep}%
  \belowdisplayskip \abovedisplayskip}
\DeclareRobustCommand\footnotesize{%
  \@setfontsize\footnotesize{\rt@pt{9}}{\rt@pt{11}}%
  \abovedisplayskip \rt@pt{8}\@plus\rt@pt{2}\@minus\rt@pt{4}%
  \abovedisplayshortskip \z@ \@plus\rt@pt{1}%
  \belowdisplayshortskip \rt@pt{4}\@plus\rt@pt{2}\@minus\rt@pt{2}%
  \def\@listi{\leftmargin\leftmargini
    \topsep \rt@pt{4}\@plus\rt@pt{2}\@minus\rt@pt{2}%
    \parsep \rt@pt{2}\@plus\rt@pt{1}\@minus\rt@pt{1}%
    \itemsep \parsep}%
  \belowdisplayskip \abovedisplayskip}
\DeclareRobustCommand\scriptsize{\@setfontsize\scriptsize{\rt@pt{8}}{\rt@pt{9.5}}}
\DeclareRobustCommand\tiny{\@setfontsize\tiny{\rt@pt{6}}{\rt@pt{7}}}
\DeclareRobustCommand\large{\@setfontsize\large{\rt@pt{12}}{\rt@pt{14}}}
\DeclareRobustCommand\Large{\@setfontsize\Large{\rt@pt{14.4}}{\rt@pt{18}}}
\DeclareRobustCommand\LARGE{\@setfontsize\LARGE{\rt@pt{17.28}}{\rt@pt{22}}}
\DeclareRobustCommand\huge{\@setfontsize\huge{\rt@pt{20.74}}{\rt@pt{25}}}
\DeclareRobustCommand\Huge{\@setfontsize\Huge{\rt@pt{24.88}}{\rt@pt{30}}}
% size11.clo's list spacing: the first level's at the normal size, which \normalsize hands to
% \@listi, and the second and third levels', the second being the bullets'.
\def\rt@listI{%
  \parsep \rt@pt{4.5}\@plus\rt@pt{2}\@minus\rt@pt{1}%
  \topsep \rt@pt{9}\@plus\rt@pt{3}\@minus\rt@pt{5}%
  \itemsep \rt@pt{4.5}\@plus\rt@pt{2}\@minus\rt@pt{1}}
\def\@listI{\leftmargin\leftmargini\rt@listI}
\def\@listii{\leftmargin\leftmarginii
  \labelwidth\leftmarginii
  \advance\labelwidth-\labelsep
  \topsep \rt@pt{4.5}\@plus\rt@pt{2}\@minus\rt@pt{1}%
  \parsep \rt@pt{2}\@plus\rt@pt{1}\@minus\rt@pt{1}%
  \itemsep \parsep}
\def\@listiii{\leftmargin\leftmarginiii
  \labelwidth\leftmarginiii
  \advance\labelwidth-\labelsep
  \topsep \rt@pt{2}\@plus\rt@pt{1}\@minus\rt@pt{1}%
  \parsep \z@
  \partopsep \rt@pt{1}\@plus\z@ \@minus\rt@pt{1}%
  \itemsep \topsep}
% \rtfontsize{<points>}: every size and its spacing from <points> in place of 11, then the normal
% size, and the first list level's spacing and \partopsep outside any list, where \trivlist, and
% so center, takes them, as size11.clo sets them as it loads. The paragraph indent stays the 0pt
% \raggedright sets.
\newcommand*\rtfontsize[1]{%
  \edef\rt@size{\number\dimexpr#1\p@\relax}%
  \normalsize
  \rt@listI
  \partopsep \rt@pt{3}\@plus\rt@pt{1}\@minus\rt@pt{1}\relax}
\makeatother
\usepackage{latexsym}
% Upstream loaded fullpage[empty] here and adjusted its margins below, which comes to half an inch
% on every side of US Letter. The page size and margins controls call \geometry again after the
% format.
\usepackage[letterpaper,margin=0.5in]{geometry}
\usepackage{titlesec}
\usepackage{marvosym}
% xcolor in place of color, and the accent colour named, as the note at the top says. The accent
% colour control defines rtaccent again after the format.
\usepackage[usenames,dvipsnames]{xcolor}
\definecolor{rtaccent}{HTML}{000000}
\usepackage{verbatim}
\usepackage{enumitem}
% dvipdfmx, because the editor makes the PDF from the DVI with xdvipdfmx: in DVI mode hyperref's
% default driver is dvips, whose link specials xdvipdfmx ignores. No bookmarks: a one-page resume
% has no use for them, and hyperref's conversion of a section title into one takes time that grows
% with the square of the title's length, about 4 seconds a pass for 8,000 characters.
\usepackage[hidelinks,dvipdfmx,bookmarks=false]{hyperref}
\usepackage{fancyhdr}
\usepackage[english]{babel}
\usepackage{tabularx}
\input{glyphtounicode}


%----------FONT OPTIONS----------
% sans-serif
% \usepackage[sfdefault]{FiraSans}
% \usepackage[sfdefault]{roboto}
% \usepackage[sfdefault]{noto-sans}
% \usepackage[default]{sourcesanspro}

% serif
% \usepackage{CormorantGaramond}
% \usepackage{charter}


\pagestyle{fancy}
\fancyhf{} % clear all header and footer fields
\fancyfoot{}
\renewcommand{\headrulewidth}{0pt}
\renewcommand{\footrulewidth}{0pt}

\urlstyle{same}

\raggedbottom
\raggedright
\setlength{\tabcolsep}{0in}

% Sections formatting
% The density control respaces what follows, after the format. Normal leaves it as it is:
% titlesec's 3.5ex before a section title and 2.3ex after it, 16.5pt and 10.8pt at 11 points, and
% the lists' topsep and itemsep, 9pt and 4.5pt for the entries and 4.5pt and 2pt for their
% bullets, which \rtfontsize above scales with the font size at the other sizes.
% Compact sets each of them lower and relaxed each higher, so that the space between sections,
% entries and bullets moves together, and both set the two levels apart: relaxed since a value
% above the entries' would more than double the bullets', and compact so that an entry stays
% further from the bullets above it than they are from each other. Compact has to leave room for
% the space the macros below take back with a negative \vspace, or it sets a line over the one
% above it, or a title's rule through the line under it. So it keeps 9pt after a title, where
% the rule's -5pt and \resumeSubheading's -2pt take back 7pt, and where a section is empty the
% rule's -5pt and the next title's -4pt take back 9pt, since titlesec drops the space before a
% title that follows another; 10pt before a title, which keeps the title a font size below an
% entry with no bullets, though the entry's heading takes back 7pt after it and the title 4pt
% before itself; and no itemsep below 0pt. At 9 points normal's space after a title, 8.9pt, is
% the one that falls under compact's, and sets nothing over anything else all the same.
\titleformat{\section}{
  \vspace{-4pt}\scshape\raggedright\large\color{rtaccent}
}{}{0em}{}[\color{rtaccent}\titlerule \vspace{-5pt}]

% Ensure that generate pdf is machine readable/ATS parsable
\pdfgentounicode=1

%-------------------------
% Custom commands
\newcommand{\resumeItem}[1]{
  \item\small{
    {#1 \vspace{-2pt}}
  }
}

\newcommand{\resumeSubheading}[4]{
  \vspace{-2pt}\item
    \begin{tabular*}{0.97\textwidth}[t]{l@{\extracolsep{\fill}}r}
      \textbf{#1} & #2 \\
      \textit{\small#3} & \textit{\small #4} \\
    \end{tabular*}\vspace{-7pt}
}

\newcommand{\resumeSubSubheading}[2]{
    \item
    \begin{tabular*}{0.97\textwidth}{l@{\extracolsep{\fill}}r}
      \textit{\small#1} & \textit{\small #2} \\
    \end{tabular*}\vspace{-7pt}
}

\newcommand{\resumeProjectHeading}[2]{
    \item
    \begin{tabular*}{0.97\textwidth}{l@{\extracolsep{\fill}}r}
      \small#1 & #2 \\
    \end{tabular*}\vspace{-7pt}
}

\newcommand{\resumeSubItem}[1]{\resumeItem{#1}\vspace{-4pt}}

\renewcommand\labelitemii{$\vcenter{\hbox{\tiny$\bullet$}}$}

\newcommand{\resumeSubHeadingListStart}{\begin{itemize}[leftmargin=0.15in, label={}]}
\newcommand{\resumeSubHeadingListEnd}{\end{itemize}}
\newcommand{\resumeItemListStart}{\begin{itemize}}
\newcommand{\resumeItemListEnd}{\end{itemize}\vspace{-5pt}}
